% Replacement prose for EXP-3; source V4 remains unchanged.
Cross-evaluating each uniform-income seat-specific design under polarized income, and each polarized-income design under uniform income, produces six infeasible candidates: every candidate exceeds the 0.95 seat-load ceiling, although finance, participation and fleet constraints are satisfied. No feasible cross-evaluated design improves the reported incumbent, so no additional search is triggered and the reported seat-pricing increments remain 395.8/169.6, 576.2/156.5 and 264.7/80.3 thousand CNY/day under uniform/polarized income in HN, WC and JX. Every retained seat-specific incumbent also exceeds its corresponding common-pricing incumbent.

The saved designs show that income composition changes which seat capacities are scarce. Moving the uniform-income seat-specific design to the polarized state raises maximum seat loads to 1.062, 1.158 and 1.057 in HN, WC and JX. In the retained polarized designs, the departure-weighted lowest G/D seat-class multipliers are lower, while the premium-seat share of H rail travelers is higher. These patterns are consistent with capacity-driven rebalancing of seat prices, but the cross-evaluations do not establish global seat-pricing optima. The reversal remains a comparison of best designs found by search, rather than a proven effect of income dispersion.
